% =============================================================================
%  revised_plan_v4.tex
%  Revision plan for "Quantum Amplitude Encoding for Discrete MRFs"
%  Author: Arul Rhik Mazumder
%  Target venue: IEEE QCE 2026, QML Track
%  Generated: 2026-04-27
%
%  This document is a planning artifact, not a submission.  It contains:
%    (1) The full set of *narrative / framing / citation* fixes I am applying
%        directly to the manuscript (no new experiments required), with the
%        corrected text inlined verbatim so it can be pasted into paper.tex.
%    (2) A precise specification of the two new classical-sampling experiments
%        that must be run before the manuscript is resubmitted:
%            E1: Exact inverse-CDF sampler  (apples-to-apples baseline)
%            E2: Parallel tempering         (stronger classical baseline)
%        Each experiment is described as: (a) protocol, (b) JSON schema for
%        the output artifact, (c) the table/figure rows it will populate,
%        (d) the exact prose that depends on its output (with TODO markers).
%    (3) A checklist mapping each reviewer concern to the fix that addresses
%        it, with a status flag (DONE / PENDING-EXPERIMENT).
%
%  Compilation: this file is meant to be read, not compiled.  It uses
%  \begin{verbatim} blocks for inlined LaTeX patches so they survive copy/paste.
% =============================================================================

\documentclass[11pt]{article}
\usepackage[margin=1in]{geometry}
\usepackage{verbatim}
\usepackage{enumitem}
\usepackage{hyperref}
\usepackage{xcolor}
\usepackage{amsmath,amssymb}
\usepackage{booktabs}

\newcommand{\DONE}{{\color{teal}\textbf{[DONE]}}}
\newcommand{\PEND}{{\color{orange}\textbf{[PENDING-EXPERIMENT]}}}
\newcommand{\TODO}[1]{{\color{red}\textbf{TODO:} #1}}

\title{Revision Plan v4 \\ \large Quantum Amplitude Encoding for Discrete MRFs}
\author{Arul Rhik Mazumder}
\date{Generated 2026-04-27}

\begin{document}
\maketitle

% =============================================================================
\section{Reviewer-Concern $\to$ Fix Map}
% =============================================================================

\begin{tabular}{p{0.55\textwidth} p{0.18\textwidth} p{0.20\textwidth}}
\toprule
\textbf{Reviewer concern} & \textbf{Fix type} & \textbf{Status} \\
\midrule
C1. Quantum advantage framing is confounded by $O(2^n)$ classical preprocessing.
& Reframing + new baseline & \PEND\ (E1 needed) \\
C2. Amplitude-encoding ``validation'' is a sanity check, not a contribution.
& Reframe \S I-D, Table XII & \DONE \\
C3. VQC results are uniformly negative; ``crossover'' framing unsupported.
& Reframe \S IV-C, \S V-A, Table XII & \DONE \\
C4. Tuned-block 1.82$\times$ is a near-tie; need a stronger classical sampler.
& New baseline & \PEND\ (E2 needed) \\
C5. Reference [7] article number is wrong (37, not 57).
& Citation fix & \DONE \\
C6. Unused references [21], [22].
& Bibliography prune & \DONE \\
C7. ``QCGM baseline'' framing is misleading (not actually QCGM).
& Reframe \S I-D, \S II, \S III-A & \DONE \\
C8. Mean-field result in Table VII deserves prominence.
& Promote to subsection & \DONE \\
C9. Abstract is dense and lacks takeaway.
& Rewrite abstract & \PEND-partial (final numbers from E1/E2) \\
C10. Compression ratios at $F<0.5$ in Table XI are misleading.
& Caveat in caption + \S IV-I prose & \DONE \\
C11. \S V-D failure mode 1 (barren plateau) contradicts \S IV-H gradient data.
& Remove failure mode 1 & \DONE \\
C12. Figure 1 is nearly empty.
& Drop fidelity panel; keep runtime panel only & \DONE \\
C13. ESS table I has no within-instance variance.
& Add seed counts + caveat & \DONE \\
C14. Table II ESS/s comparison is hardware-incomparable.
& Caveat + retitle & \DONE \\
C15. No hardware experiments.
& Acknowledge as a stated scope limit & \DONE \\
\bottomrule
\end{tabular}

\medskip
\noindent\textbf{Net status:} 12 of 15 concerns resolved by text edits in this plan.
The remaining 3 (C1, C4, and the final numeric values in C9) require running
experiments E1 and E2 below.

% =============================================================================
\section{New Experiments to Run (E1, E2)}
% =============================================================================

\subsection{E1 --- Exact Inverse-CDF Baseline (apples-to-apples)}

\paragraph{Motivation.}  The amplitude-encoding pipeline already computes the
full diagonal $H_\theta$ in $O(2^n)$ time and normalizes it to $P_\theta(x)$.
Once that array is in memory, the correct classical baseline is exact
inverse-CDF sampling, not Gibbs --- both methods amortize the same
preprocessing cost and both produce i.i.d.\ samples with $\tau \approx 1$.
This baseline directly addresses reviewer concern C1 and is the most
important single addition to the paper.

\paragraph{Protocol.}
\begin{itemize}[leftmargin=*]
  \item For every instance already in the 60-instance ESS benchmark
        (5 families $\times$ 12 instances), reuse the exact same
        precomputed $P_\theta$ array used by the quantum sampler.
  \item Build the cumulative array $F(j) = \sum_{i \le j} P_\theta(x_i)$
        once per instance ($O(2^n)$ time, one-time cost).
  \item Draw 3000 i.i.d.\ samples per instance via binary search on $F$
        (each sample is $O(n)$).  No burn-in.
  \item Record per-instance:
        \begin{itemize}
          \item \verb|t_preprocess_s|: time to build the cumulative array
          \item \verb|t_sample_s|: wall-clock for the 3000 samples (excludes preprocess)
          \item \verb|t_total_s|: \verb|t_preprocess_s| + \verb|t_sample_s|
          \item \verb|ess|: from the same scalar series + autocorrelation estimator used elsewhere
                (should be $\approx 3000$; report it explicitly --- do not assume)
          \item \verb|ess_per_s_sample_only|: \verb|ess / t_sample_s|
          \item \verb|ess_per_s_amortized|: \verb|ess / t_total_s|
        \end{itemize}
  \item For the quantum sampler row in the comparison table, also report a
        new \verb|ess_per_s_amortized| number that includes the
        $O(2^n)$ Hamiltonian-diagonal preprocessing cost.  This is the
        key fairness fix.
\end{itemize}

\paragraph{Output artifact.}  Write to
\verb|experiments/results/inverse_cdf_baseline.json| with the schema below.
This file plus an updated \verb|paper_table_metrics.json| are the only new
artifacts E1 produces.

\begin{verbatim}
{
  "generated_utc": "<ISO8601>",
  "n_instances": 60,
  "families": ["barbell", "barbell_path", "chain",
               "erdos_renyi", "two_clique"],
  "burn_in": 0,
  "samples_per_instance": 3000,
  "per_instance": [
    {
      "instance_id": "<string>",
      "family": "<family>",
      "n": <int>,
      "t_preprocess_s": <float>,
      "t_sample_s": <float>,
      "t_total_s": <float>,
      "ess": <float>,
      "ess_per_s_sample_only": <float>,
      "ess_per_s_amortized": <float>
    }, ...
  ],
  "summary": {
    "mean_ess": <float>,
    "median_ess": <float>,
    "mean_ess_per_s_sample_only": <float>,
    "mean_ess_per_s_amortized": <float>
  }
}
\end{verbatim}

\paragraph{What this populates.}  A new row in the (renamed) wall-clock
diagnostic table, plus replacement prose in \S IV-A and the abstract.
Exact text is given in \S\ref{sec:revised-text} with TODO markers for the
numeric placeholders.

\subsection{E2 --- Parallel Tempering Baseline (stronger classical)}

\paragraph{Motivation.}  The current strongest classical baseline is
tuned-block Gibbs (mean ratio $1.82\times$).  Reviewer concern C4 notes this
is a near-tie and that modern MCMC for peaked discrete distributions
typically uses tempering.  Adding a parallel-tempering (PT) baseline tests
whether the residual quantum advantage survives against a state-of-the-art
discrete sampler.

\paragraph{Protocol.}
\begin{itemize}[leftmargin=*]
  \item For each of the 60 instances:
    \begin{itemize}
      \item Run $K = 8$ replicas at inverse temperatures
            $\beta_k = \beta_{\min} + (\beta_{\max} - \beta_{\min})\,k/(K-1)$
            with $\beta_{\min} = 0.1$, $\beta_{\max} = 1.0$ (target distribution
            corresponds to $\beta = 1$).  The geometric ladder
            $\beta_k = \beta_{\min}(\beta_{\max}/\beta_{\min})^{k/(K-1)}$ is
            also acceptable; report whichever is used.
      \item Each replica runs single-site Gibbs internally (so this is a
            ``parallel-tempered single-site'' baseline).
      \item Replica swap proposals every $L_{\text{swap}} = 10$ within-replica
            sweeps, using the standard Metropolis swap acceptance ratio.
      \item Burn-in: 1000 within-replica sweeps per replica, matching the
            other Gibbs baselines.
      \item Retain 3000 samples from the $\beta = 1$ replica only.
      \item Tune $\beta_{\min}$ once on a held-out instance per family to target
            $\sim$30\% swap acceptance between adjacent rungs; do not re-tune
            per-instance (keeps the comparison honest).
    \end{itemize}
  \item Record per-instance:
        \verb|ess|, \verb|ess_per_s|, \verb|swap_acceptance_mean|,
        \verb|t_total_s| (sum over all replicas, since PT pays for all of them).
  \item Compute Quantum/PT ESS ratios over the 60 instances; report
        mean / median / range exactly as for the other three Gibbs comparators.
\end{itemize}

\paragraph{Output artifact.}  Write to
\verb|experiments/results/parallel_tempering_baseline.json|:

\begin{verbatim}
{
  "generated_utc": "<ISO8601>",
  "n_instances": 60,
  "n_replicas": 8,
  "beta_schedule": "geometric",   // or "linear"
  "beta_min": 0.1,
  "beta_max": 1.0,
  "swap_interval_sweeps": 10,
  "burn_in_per_replica_sweeps": 1000,
  "samples_per_instance": 3000,
  "per_instance": [
    {
      "instance_id": "<string>",
      "family": "<family>",
      "n": <int>,
      "ess": <float>,
      "t_total_s": <float>,
      "ess_per_s": <float>,
      "swap_acceptance_mean": <float>
    }, ...
  ],
  "summary": {
    "ratio_quantum_over_pt_mean": <float>,
    "ratio_quantum_over_pt_median": <float>,
    "ratio_quantum_over_pt_min": <float>,
    "ratio_quantum_over_pt_max": <float>,
    "mean_ess_per_s": <float>,
    "mean_swap_acceptance": <float>
  }
}
\end{verbatim}

\paragraph{Pre-registered prediction.}  I expect the Quantum/PT ratio to be
in the range $[0.8, 1.5]$ at the mean, i.e., PT will close the gap further
than tuned-block did and may even invert it on some instances.  This is the
honest expected outcome and the paper text below is written to handle that
case gracefully.  If the ratio comes in $> 2$, that would itself be a notable
finding worth foregrounding; if $< 1$ on average, the paper should explicitly
concede that the quantum i.i.d.\ advantage does not survive against modern
classical samplers in this regime.

% =============================================================================
\section{Revised Manuscript Text (paste-ready)}
\label{sec:revised-text}
% =============================================================================

The remainder of this document contains LaTeX patches.  Each patch is a
\verb|verbatim| block containing the exact replacement text for a labeled
section of \texttt{paper.tex}.  Apply them in order.  TODO markers indicate
the few places where E1/E2 numeric outputs must be inserted before
resubmission.

% -----------------------------------------------------------------------------
\subsection*{Patch 1 --- Title}
% -----------------------------------------------------------------------------

\textbf{Old:} ``Quantum Amplitude Encoding for Discrete MRFs: An ESS
Characterization Against Single-Site Gibbs Baselines.''

\noindent\textbf{New:}

\begin{verbatim}
\title{Amplitude-Encoded Sampling for Discrete MRFs:
       A Characterization of Autocorrelation Costs Across
       Classical Baselines}
\end{verbatim}

\noindent\emph{Rationale:} The previous title singled out single-site Gibbs,
which is the weakest of four baselines.  The new title is honest about scope:
this is a characterization of autocorrelation costs, not a quantum-advantage
claim.

% -----------------------------------------------------------------------------
\subsection*{Patch 2 --- Abstract (replaces existing abstract entirely)}
% -----------------------------------------------------------------------------

\begin{verbatim}
\begin{abstract}
Sampling from discrete Markov random fields (MRFs) is a well-known
hard problem in probabilistic inference.  We study amplitude-encoded
i.i.d. sampling for small discrete MRFs and characterize the
autocorrelation cost of four classical sampling baselines on matched
instances.  Our setup is deliberately scoped to the regime where the
$2^n$ probabilities of the target distribution can be precomputed
classically; in this regime no quantum exponential speedup is possible,
but the structural property of the amplitude-encoded sampler --- that
each circuit execution returns an independent sample with effective
sample size $\tau \approx 1$ --- can be cleanly compared against
classical alternatives.

Across 60 instances spanning five graph families (barbell, barbell-path,
chain, Erdos--Renyi, two-clique) with 1\,000-step burn-in and 3\,000
retained samples, Quantum/Single-Site-Gibbs ESS ratios have mean 16.35
(median 15.99, range $[8.24, 30.31]$); Quantum/Block-Gibbs ratios have
mean 7.29 (median 6.78, range $[3.31, 18.30]$); Quantum/Tuned-Block-Gibbs
ratios have mean 1.82 (median 1.91, range $[1.03, 2.43]$); and
Quantum/Parallel-Tempering ratios have mean TODO-E2-MEAN
(median TODO-E2-MEDIAN, range TODO-E2-RANGE).  This hierarchy quantifies
the autocorrelation cost of each MCMC design choice and shows that
modern classical samplers substantially close, and may eliminate, the
ESS gap relative to amplitude-encoded i.i.d.\ sampling.

Critically, when the $O(2^n)$ classical preprocessing required by the
amplitude encoder is amortized into the wall-clock comparison, exact
inverse-CDF sampling on the precomputed distribution achieves
TODO-E1-ESS-PER-S-AMORTIZED ESS/s versus
TODO-QUANTUM-ESS-PER-S-AMORTIZED for the quantum sampler --- confirming
that the quantum sampler offers no wall-clock advantage in the regime
studied.  The contribution of this work is therefore not a speedup
claim but (i) a clean characterization of MCMC autocorrelation costs
under a fixed protocol, and (ii) a reproducible benchmark of
amplitude-encoded state preparation for discrete MRFs at $n=8,10,12$
on simulator backends.

We also report a multi-trial MPS scaling study (three seeds per point,
$n$ up to 40) showing $\chi=32$ achieves $F = 0.721 \pm 0.059$ at
$n = 40$, providing a concrete classical benchmark ceiling.  A
matched-budget VQC vs.\ MPS comparison at $n = 8, 10, 12$ shows VQC
fidelities below MPS at every measured point
($(F_{\text{VQC}}, F_{\text{MPS}}) = (0.306, 0.990), (0.210, 0.958),
(0.165, 0.878)$ at parameter compressions of $10.7\times$, $34.1\times$,
$113.8\times$); we report this as a negative result for shallow
hardware-efficient ansatze in this regime rather than a compression
finding.  All code and data are at
https://github.com/arulrhikm/QuantumDGM.
\end{abstract}
\end{verbatim}

\noindent\emph{Rationale:} The new abstract (i) leads with scope
limitations, (ii) reports the PT comparator alongside the others, (iii)
explicitly amortizes preprocessing into the wall-clock comparison, (iv)
reframes VQC as a negative result, and (v) drops the speedup framing
entirely.

% -----------------------------------------------------------------------------
\subsection*{Patch 3 --- \S I-A (Problem Selection and Motivation),
last paragraph replacement}
% -----------------------------------------------------------------------------

\textbf{Old (last paragraph):} ``These limitations together motivate the
exploration of quantum approaches that, in the amplitude-encoding path, can
produce exact samples from the target distribution without any notion of
mixing time.''

\noindent\textbf{New:}

\begin{verbatim}
These limitations of MCMC --- slow mixing, burn-in cost, sample
correlation --- have motivated quantum approaches.  We emphasize at
the outset that the amplitude-encoding path studied here does
\emph{not} circumvent the $O(2^n)$ cost of computing the target
distribution; it requires building the full diagonal of the MRF
Hamiltonian classically before any quantum operation.  In this regime
the quantum sampler is best understood as an alternative i.i.d.\
sample generator that runs after the classical preprocessing has
already enumerated $P_\theta$.  The relevant scientific questions are
therefore (i) how much autocorrelation cost is paid by classical
MCMC alternatives that share the same preprocessing budget, and
(ii) whether any structural property of the quantum sampler beyond
i.i.d.\ output (e.g., zero burn-in, parallelism) translates into a
practical advantage.  We address (i) directly with the ESS
characterization in \S IV-A and address (ii) honestly --- including
where the answer is negative --- in \S V.
\end{verbatim}

% -----------------------------------------------------------------------------
\subsection*{Patch 4 --- \S I-D (Baseline and Novel Contributions),
contributions list replacement}
% -----------------------------------------------------------------------------

\textbf{Old:} The four-item contribution list and the framing of the
amplitude encoder as a ``QCGM-based baseline.''

\noindent\textbf{New:}

\begin{verbatim}
We construct an amplitude-encoding pipeline for discrete MRFs that
prepares $\ket{\psi} = \sum_x \sqrt{P_\theta(x)}\,\ket{x}$ via Qiskit's
\verb|StatePreparation| primitive after classical computation of the
target diagonal.  We emphasize that this is \emph{not} the QCGM circuit
construction of Piatkowski and Zoufal~\cite{piatkowski2024qcgm}, which
uses ancilla-based real-part extraction and a repeat-until-success
sampling scheme with a unitary Hammersley--Clifford structure.  Our
pipeline is a simpler, more directly auditable amplitude encoder that
sacrifices the potential exponential speedup of full QCGM Hamiltonian
simulation for clean validation of the small-$n$ sampling regime.

With that scope fixed, our contributions are:
\begin{enumerate}
\item A characterization of the ESS gap between amplitude-encoded
      i.i.d.\ sampling and four classical samplers (single-site Gibbs,
      block Gibbs, tuned-block Gibbs, and parallel tempering) across
      60 instances stratified by graph family.  This is the primary
      finding of the paper.
\item A wall-clock comparison that explicitly amortizes the $O(2^n)$
      preprocessing cost on both sides, including a comparison against
      exact inverse-CDF sampling --- the natural apples-to-apples
      classical baseline once the full distribution has been
      enumerated.
\item A multi-trial MPS scaling curve through $n = 40$ with three seeds
      per point, providing a concrete classical benchmark ceiling for
      future variational and quantum approaches at large $n$.
\item Reproducibility infrastructure: an open-source implementation
      with executable artifact-verification scripts that regenerate
      every table and figure from committed JSON files.
\end{enumerate}
We also report several negative or null results that we believe are
useful to the community: shallow hardware-efficient VQC ansatze
underperform matched-budget MPS at every measured size $n \le 12$ with
no evidence of a forthcoming crossover; clique-aware entanglement
provides no measurable benefit over linear entanglement at tested
scales; and amplitude-encoded VQC is beaten by mean-field on full
distributional fidelity at $n = 8$.
\end{verbatim}

\noindent\emph{Rationale:} Drops the amplitude-encoding ``benchmark'' as
a contribution (it is a sanity check on a Qiskit primitive --- C2);
explicitly disclaims the QCGM framing (C7); replaces VQC as a contribution
with VQC as a negative result (C3); and adds the inverse-CDF and PT
baselines (C1, C4).

% -----------------------------------------------------------------------------
\subsection*{Patch 5 --- \S III-A (Amplitude Encoding), ``Design rationale''
paragraph}
% -----------------------------------------------------------------------------

\textbf{Old:} ``We chose the simplified amplitude-encoding variant over the
full QCGM circuit construction because it is easier to audit/debug and
already supports strong small-$n$ fidelity in the current artifacts.''

\noindent\textbf{New:}

\begin{verbatim}
\paragraph{Design rationale and scope.} Our pipeline is a
\emph{simplified amplitude encoder}, not a QCGM circuit.  The QCGM
construction~\cite{piatkowski2024qcgm} embeds the graphical model
unitarily and uses ancilla-based real-part extraction in a
repeat-until-success scheme, achieving a sampling success probability
that does not depend on the number of variables.  By contrast, we
classically precompute the diagonal of $H_\theta$, normalize to obtain
$P_\theta$, and call \texttt{StatePreparation} on the resulting
amplitudes.  This trades the potential exponential speedup of QCGM
Hamiltonian simulation for two practical benefits in the small-$n$
regime: (i) the entire pipeline can be audited end-to-end against
direct enumeration of $P_\theta$, and (ii) any deviation from
$F = 1$ is attributable to a small number of well-localized causes
(bit-ordering, normalization, backend numerics).  We emphasize that
the $O(2^n)$ classical preprocessing this approach requires is the
same cost paid by any classical method that operates on the
fully-enumerated distribution, including exact inverse-CDF sampling.
\end{verbatim}

% -----------------------------------------------------------------------------
\subsection*{Patch 6 --- \S III-D (Experimental Setup),
``ESS baseline protocol'' paragraph extension}
% -----------------------------------------------------------------------------

\textbf{Append after the existing ``ESS baseline protocol'' paragraph:}

\begin{verbatim}
\paragraph{Additional baselines (E1, E2).}
We add two baselines beyond the three Gibbs variants:

\begin{itemize}
\item \textbf{Exact inverse-CDF sampling.} For each instance, we
      reuse the precomputed $P_\theta$ array, build the cumulative
      distribution $F(j) = \sum_{i \le j} P_\theta(x_i)$, and draw
      3\,000 i.i.d.\ samples via binary search.  This is the
      apples-to-apples classical baseline once the full distribution
      has been enumerated.  We report two wall-clock costs: a
      sample-only cost that excludes preprocessing (matching the
      quantum sampler's per-shot cost) and an amortized cost that
      includes the one-time $O(2^n)$ cumulative-array construction
      (matching the quantum sampler's full pipeline).

\item \textbf{Parallel tempering (PT).} A standard PT sampler with
      $K=8$ replicas, geometric inverse-temperature ladder
      $\beta \in [\beta_{\min}, 1]$ tuned per family for $\sim$30\%
      swap acceptance, replica swap proposals every 10 within-replica
      single-site Gibbs sweeps, 1\,000 burn-in sweeps per replica,
      and 3\,000 retained samples from the $\beta = 1$ replica.  PT
      total wall-clock includes all replica work.
\end{itemize}

These two additions tighten the comparison in opposite directions:
inverse-CDF establishes the ``cheapest possible'' classical baseline
that shares the quantum preprocessing cost, while PT establishes a
``stronger than tuned-block'' classical sampler that does not require
full enumeration of $P_\theta$.
\end{verbatim}

% -----------------------------------------------------------------------------
\subsection*{Patch 7 --- \S III-D, replace the ``Hardware'' paragraph}
% -----------------------------------------------------------------------------

\textbf{Old:} The existing ``Hardware'' paragraph.

\noindent\textbf{New:}

\begin{verbatim}
\paragraph{Backends and scope of hardware claims.}
All quantum experiments use simulator backends: local Aer for
small-$n$ verification and BlueQubit cloud statevector / MPS
simulators for the published artifact runs.  No physical quantum
hardware results are reported in this submission, and none of the
ESS, fidelity, or compression numbers should be interpreted as
predictive of noisy-hardware performance.  Wall-clock timings on
BlueQubit are subject to queue latency and should be treated as
order-of-magnitude indicators rather than precise benchmarks; we
discuss this caveat further alongside Table~\ref{tab:wallclock}.
\end{verbatim}

\noindent\emph{Rationale:} Acknowledges C15 (no hardware) up front rather
than deferring to \S V.

% -----------------------------------------------------------------------------
\subsection*{Patch 8 --- \S IV opening paragraph}
% -----------------------------------------------------------------------------

\textbf{Old:} ``Our results section is organized to lead with the primary
completed finding (ESS characterization, \S IV-A), followed by amplitude
encoding correctness (\S IV-B), the variational crossover experiment
(\S IV-C), and MPS/amplitude scaling in progress (\S IV-I).''

\noindent\textbf{New:}

\begin{verbatim}
The results are organized as follows.  \S IV-A reports the primary
finding: ESS characterization across four classical baselines
including parallel tempering, plus the wall-clock comparison against
exact inverse-CDF sampling.  \S IV-B is a sanity check on the
amplitude-encoding pipeline (it is a verification artifact, not a
research finding).  \S IV-C reports the VQC-vs-MPS matched-budget
comparison as a negative result for shallow ansatze in this regime.
\S IV-D and \S IV-E document the entanglement-strategy null result.
\S IV-F highlights an instructive negative result: VQC is beaten by
mean-field on full distributional fidelity at $n = 8$.  \S IV-G
analyzes distribution hardness; \S IV-H reports VQC convergence and
gradient-norm diagnostics; \S IV-I is the multi-trial MPS scaling
study through $n = 40$.
\end{verbatim}

% -----------------------------------------------------------------------------
\subsection*{Patch 9 --- \S IV-A (Primary Finding), full replacement}
% -----------------------------------------------------------------------------

\begin{verbatim}
\subsection{ESS Characterization Across Four Classical Baselines}
\label{sec:ess}

Quantum amplitude encoding produces samples that are structurally
independent across circuit executions, yielding $\tau \approx 1$
regardless of graph topology.  We quantify the resulting ESS gap
against four classical baselines (single-site Gibbs, block Gibbs,
tuned-block Gibbs, and parallel tempering) over 60 matched instances
under the protocol in \S III-D.

Results are summarized in Table~\ref{tab:ess-summary} and
Fig.~\ref{fig:ess-bars}.  The four ratios show a clear monotone
hierarchy:
\begin{itemize}
\item Quantum/Single-Site-Gibbs: mean 16.35, median 15.99, range
      $[8.24, 30.31]$.
\item Quantum/Block-Gibbs: mean 7.29, median 6.78, range
      $[3.31, 18.30]$.
\item Quantum/Tuned-Block-Gibbs: mean 1.82, median 1.91, range
      $[1.03, 2.43]$.
\item Quantum/Parallel-Tempering: mean TODO-E2-MEAN, median
      TODO-E2-MEDIAN, range TODO-E2-RANGE.
\end{itemize}

The interpretation depends on the PT result.  If the PT mean ratio
is greater than approximately 1.5, the structural i.i.d.\ property of
the amplitude-encoded sampler still translates into a measurable
autocorrelation advantage even against a modern classical sampler.
If the PT mean ratio is at or below 1, the i.i.d.\ property does not
translate into a practical autocorrelation advantage in this regime,
and the quantum sampler's value lies elsewhere (e.g., in the absence
of burn-in or in suitability for downstream parallel sampling
workflows).  We report whichever conclusion the data supports.
TODO-E2-INTERPRETATION-PARAGRAPH.

\paragraph{Wall-clock with amortized preprocessing.}
The above ESS ratios characterize autocorrelation cost in
sample-count units.  To assess whether the structural advantage
translates into a wall-clock advantage, we compare per-instance
ESS/s including the $O(2^n)$ classical preprocessing on both sides
(Table~\ref{tab:wallclock}).  Exact inverse-CDF sampling on the
precomputed distribution achieves TODO-E1-ESS-PER-S-AMORTIZED ESS/s
amortized, versus TODO-QUANTUM-ESS-PER-S-AMORTIZED for the quantum
sampler.  TODO-E1-INTERPRETATION (one of: ``Inverse-CDF dominates the
amortized comparison by a factor of TODO-X, confirming that the
amplitude-encoded sampler offers no wall-clock advantage in this
regime'' / ``The two are within a factor of TODO-X, indicating that
the quantum sampler is wall-clock competitive in this regime once
preprocessing is amortized'').  In either case, the relevant
takeaway is that any practical advantage of the amplitude-encoded
sampler must come from properties other than raw ESS/s --- for
example, the absence of any burn-in requirement, exact i.i.d.\
output for statistical testing, or natural parallelism across
circuit executions.

\paragraph{Topology effects dominate a single spectral-gap
predictor.} The ESS-ratio versus mixing-difficulty relationship
(Fig.~\ref{fig:ess-mixing}) shows a weak global linear trend
($R^2 \approx 0.00036$ for the single-site comparator), confirming
that family/topology effects dominate over a single spectral-gap
predictor in our 60-instance benchmark.  This is a property of the
benchmark, not a finding about quantum versus classical sampling per
se: future work should evaluate on instance distributions with
larger spread in mixing difficulty (e.g., low-temperature spin
glasses with $\theta$ scaled to push the distribution further into
the combinatorial-optimization regime).
\end{verbatim}

% -----------------------------------------------------------------------------
\subsection*{Patch 10 --- Table I (replace) and Table II (replace + retitle)}
% -----------------------------------------------------------------------------

\begin{verbatim}
\begin{table}[t]
\caption{ESS-ratio summary across four classical baselines (60
instances, 5 graph families, 1\,000-step burn-in, 3\,000 retained
samples per instance).  All values from
\texttt{paper\_table\_metrics.json} generated 2026-XX-XX after
incorporating E1 and E2.  ESS values within an instance are point
estimates from a single chain via integrated autocorrelation; we
caution that these per-instance estimates are noisy and the
ratios should be interpreted at the distributional level (mean,
median, range) rather than per-instance.}
\label{tab:ess-summary}
\centering
\begin{tabular}{lrrrr}
\toprule
Comparator & Mean & Median & Min & Max \\
\midrule
Quantum / Single-site Gibbs        & 16.35  & 15.99  & 8.24  & 30.31 \\
Quantum / Block Gibbs              & 7.29   & 6.78   & 3.31  & 18.30 \\
Quantum / Tuned-Block Gibbs        & 1.82   & 1.91   & 1.03  & 2.43  \\
Quantum / Parallel Tempering       & TODO-PT-MEAN & TODO-PT-MED
                                   & TODO-PT-MIN  & TODO-PT-MAX \\
\bottomrule
\end{tabular}
\end{table}
\end{verbatim}

\begin{verbatim}
\begin{table}[t]
\caption{Wall-clock-amortized ESS/s diagnostics, including
preprocessing cost on both sides (mean over 60 instances).
``Sample-only'' columns exclude the $O(2^n)$ preprocessing;
``Amortized'' columns include it.  The amortized comparison is the
fair one when the full distribution must be enumerated.  Quantum
runs use BlueQubit SV simulator; classical runs use local CPU.
We caution that backend heterogeneity (cloud vs.\ local, queue
latency, simulator implementation) makes absolute numbers
order-of-magnitude indicative rather than precise.}
\label{tab:wallclock}
\centering
\begin{tabular}{lrr}
\toprule
Method & Sample-only ESS/s & Amortized ESS/s \\
\midrule
Quantum amplitude-encoded sampler & 714.01
                                  & TODO-Q-AMORT \\
Exact inverse-CDF                 & TODO-E1-SAMP
                                  & TODO-E1-AMORT \\
Single-site Gibbs                 & 3.43  & 3.43  \\
Block Gibbs                       & 0.089 & 0.089 \\
Tuned-block Gibbs                 & 0.120 & 0.120 \\
Parallel tempering                & TODO-PT-ESS-S & TODO-PT-ESS-S \\
\bottomrule
\end{tabular}
\end{table}
\end{verbatim}

\noindent\emph{Rationale:} Table I now reports four baselines (C4); Table II
explicitly amortizes preprocessing (C1) and adds the backend-heterogeneity
caveat (C14).  Note that Gibbs and PT do not require full enumeration, so
their sample-only and amortized columns are equal.

% -----------------------------------------------------------------------------
\subsection*{Patch 11 --- \S IV-B (Amplitude Encoding), reframe as
sanity check}
% -----------------------------------------------------------------------------

\textbf{Old:} ``Table III and Fig. 1 report generated fidelity values for
amplitude encoding on chain graphs of varying size [...] validating the
Hamiltonian normalization, bit-ordering convention, and diagonal-only memory
optimization end-to-end.''

\noindent\textbf{New:}

\begin{verbatim}
\subsection{Amplitude-Encoding Pipeline Verification}
\label{sec:amp-verify}

Table~\ref{tab:amp} reports a sanity-check verification of the
amplitude-encoding pipeline on chain graphs at $n = 8, 10, 12$.  All
three runs achieve $F \approx 1.000$ and $\mathrm{TV} = 0.000$ on
BlueQubit statevector backends with no Aer fallback.  This is the
expected outcome of a correct call to \texttt{StatePreparation} on a
statevector simulator and we present it as a verification artifact
that the bit-ordering convention, normalization, and diagonal-only
memory optimization are implemented correctly --- not as a research
finding.  Any failure to achieve $F \approx 1$ here would indicate a
bug in the pipeline; achieving it confirms correctness but is not
itself novel.
\end{verbatim}

\noindent\emph{Rationale:} Addresses C2 directly.  Also: drop Fig.\ 1's
fidelity panel (C12) since three points at $F = 1$ is uninformative; keep
the runtime panel only and merge it inline with Table~\ref{tab:amp}.

% -----------------------------------------------------------------------------
\subsection*{Patch 12 --- \S IV-C (Variational Compression),
full reframing}
% -----------------------------------------------------------------------------

\textbf{Old:} ``Table V summarizes the currently generated Experiment B
artifact rows [...] we observe the expected compression/accuracy trade-off
under matched-budget comparison, with MPS reaching higher fidelity than the
current shallow VQC runs at these sizes.''

\noindent\textbf{New:}

\begin{verbatim}
\subsection{Shallow VQC Underperforms MPS: A Negative Result}
\label{sec:vqc-neg}

Table~\ref{tab:largeN} reports matched-budget comparisons of a
depth-3 hardware-efficient VQC against MPS at $\chi_{\max} = 2$
across $n = 8, 10, 12$.  At every measured size the VQC fidelity is
substantially below MPS:
$(F_{\text{VQC}}, F_{\text{MPS}}) = (0.306, 0.990)$ at $n = 8$,
$(0.210, 0.958)$ at $n = 10$, $(0.165, 0.878)$ at $n = 12$.

We frame this as a negative result for shallow hardware-efficient
ansatze in the peaked-MRF regime, not as a compression-vs-accuracy
trade-off.  Although the VQC achieves nominal parameter compressions
of $10.7\times$, $34.1\times$, $113.8\times$, the absolute fidelities
are below the threshold at which the resulting samples would be
useful for any downstream marginal- or moment-estimation task; a
distribution with $F = 0.165$ overlaps the target by 16.5\%, and
samples drawn from it would not match $P_\theta$ in any practical
sense.  The fidelity gap narrows from $-0.748$ at $n = 10$ to
$-0.713$ at $n = 12$ (a $0.035$ decrease over 2 qubits).  We
acknowledge that this rate of narrowing does not support
extrapolating a near-term VQC--MPS crossover, and we do not claim
one.  Discovering whether a positive crossover exists at any
larger $n$ for this ansatz family is left as open future work.

\paragraph{What this means for the variational path.}  Two natural
explanations for the observed fidelity gap are (i) insufficient
circuit depth at these problem sizes and (ii) the limited
expressivity of the depth-3 hardware-efficient ansatz relative to
the highly peaked MRF distributions documented in
Table~\ref{tab:hardness}.  The convergence and gradient-norm
diagnostics in \S IV-H rule out a simple barren-plateau or
optimizer-divergence explanation, leaving expressivity as the
most likely binding constraint.
\end{verbatim}

\noindent\emph{Rationale:} Addresses C3 directly --- drops the ``crossover
may exist at $n^* > 12$'' claim that the data does not support.

% -----------------------------------------------------------------------------
\subsection*{Patch 13 --- \S IV-F (Classical Baseline Comparison), promote
mean-field finding}
% -----------------------------------------------------------------------------

\textbf{Old:} ``Table VII compares quantum amplitude encoding against two
standard classical approximation methods [...]''

\noindent\textbf{New:}

\begin{verbatim}
\subsection{Mean-Field Beats VQC on Full Distributional Fidelity}
\label{sec:meanfield}

Table~\ref{tab:baselines} compares the depth-3 VQC against two
standard classical approximation methods at $n = 6, 8, 10$.  The
notable finding is that mean-field approximation \emph{beats} the
VQC on full distributional fidelity at $n = 8$ (0.787 vs.\ 0.664)
and $n = 10$ (0.606 vs.\ 0.556).  This is an honest negative result
for the variational path: a textbook tractable approximation
outperforms a parameterized quantum circuit with $O(nd)$ trainable
parameters in matching the global distribution.

The picture is more nuanced for marginal estimation: the VQC
achieves \emph{lower} singleton-marginal mean absolute error than
both mean-field and loopy belief propagation at every measured
size.  This indicates that the depth-3 ansatz captures
\emph{individual-variable marginals} more accurately than these
classical approximations even when its global fit is worse.  In
practice this means the VQC may be a useful approximation if the
downstream task is marginal estimation but is not a useful
approximation if the downstream task requires faithful joint
samples.  Amplitude encoding (which trivially attains $F = 1.000$
on this benchmark) is shown for reference and is not a member of
the baseline set.
\end{verbatim}

\noindent\emph{Rationale:} Addresses C8 (promote mean-field finding to a
labeled subsection).

% -----------------------------------------------------------------------------
\subsection*{Patch 14 --- \S IV-I caption + Table XI caption update}
% -----------------------------------------------------------------------------

\textbf{Add to the end of the Table XI caption:}

\begin{verbatim}
Caveat on compression at low fidelity: rows with $F < 0.5$ (e.g.,
$\chi = 8$ at $n = 32, 40$ and $\chi = 16$ at $n = 32, 40$) are
reported for completeness but the corresponding compression ratios
should not be interpreted as evidence of useful approximation: at
$F < 0.5$ the MPS distribution is closer to uniform than to
$P_\theta$ on most observables.  The headline scaling result is the
$\chi = 32$ row at $n = 40$ ($F = 0.721 \pm 0.059$), which is the
only point at $n \ge 24$ that exceeds $F = 0.7$.
\end{verbatim}

\noindent\emph{Rationale:} Addresses C10.

% -----------------------------------------------------------------------------
\subsection*{Patch 15 --- \S V-A (Key Findings and Insights), full rewrite}
% -----------------------------------------------------------------------------

\begin{verbatim}
\subsection{Key Findings and Insights}

Five findings emerge from the committed artifacts, ordered by
strength of the conclusion they support:

\begin{itemize}
\item \textbf{Autocorrelation hierarchy across four baselines.}
The Quantum/baseline ESS ratios are 16.35 (single-site Gibbs),
7.29 (block Gibbs), 1.82 (tuned-block Gibbs), and TODO-E2-MEAN
(parallel tempering).  The progression from a $16\times$ ratio
against the weakest baseline to a near-tie or below-unity ratio
against the strongest is the most informative result of this
work.  It quantifies, on a single instance set, exactly how much
of the apparent ``quantum sampling advantage'' is actually
attributable to choosing a weak classical baseline rather than to
the quantum sampler itself.

\item \textbf{No wall-clock advantage in the precomputed-$P_\theta$
regime.} When the $O(2^n)$ classical preprocessing required by the
amplitude encoder is amortized into the comparison, exact
inverse-CDF sampling achieves TODO-E1-RATIO\,$\times$ the amortized
ESS/s of the quantum sampler.  The structural i.i.d.\ property of
the amplitude encoder does not translate into a wall-clock advantage
in this regime; any practical case for amplitude encoding must rest
on properties other than throughput (e.g., absence of burn-in,
exact i.i.d.\ output for testing, parallelism).

\item \textbf{MPS scales to $n = 40$ with confidence intervals.}
Three-seed evaluation confirms $\chi = 32$ achieves
$F = 0.721 \pm 0.059$ at $n = 40$, establishing a concrete
classical performance ceiling for future variational and quantum
approaches at large $n$.

\item \textbf{Shallow VQC underperforms MPS at every measured size.}
At $n = 8, 10, 12$, depth-3 VQC fidelities are 3--6$\times$ below
matched-budget MPS, and the gap-narrowing rate does not support
extrapolating a near-term crossover.  This is an honest negative
result for hardware-efficient ansatze on peaked MRF distributions.

\item \textbf{Mean-field beats VQC on global fit; VQC beats both
mean-field and loopy BP on singleton marginals.}  This separation
suggests the variational quantum approach studied here is a
candidate for marginal-estimation workflows but not for
joint-sample workflows, in the regime tested.
\end{itemize}
\end{verbatim}

% -----------------------------------------------------------------------------
\subsection*{Patch 16 --- \S V-D (Limitations and Failure Modes), edits}
% -----------------------------------------------------------------------------

\textbf{Edits to the existing limitation list:}

\begin{itemize}[leftmargin=*]
\item \textbf{Replace} the first bullet (``No exponential speedup'') with a
      shorter, sharper version that matches the paper's new framing:

\begin{verbatim}
\item \textbf{Scope of the studied regime.} The amplitude-encoding
pipeline studied here requires $O(2^n)$ classical preprocessing and
is therefore restricted to small-$n$ settings where this cost is
acceptable.  This is the same regime in which exact classical
samplers (inverse-CDF after enumeration) are feasible.  We make no
claims about quantum performance outside this regime; in
particular, we do not implement the QCGM construction of
Piatkowski and Zoufal that could in principle achieve an
exponential speedup via Hamiltonian simulation, nor do we evaluate
on noisy hardware.
\end{verbatim}

\item \textbf{Delete} ``Failure mode 1: Barren plateaus.''  The
      gradient-norm data in \S IV-H rules out a barren-plateau
      explanation; including it as a failure mode at $n = 12$
      contradicts our own measurements (C11).

\item \textbf{Edit} the ``What the ESS comparison does and does not show''
      paragraph to add the PT comparator and inverse-CDF context:

\begin{verbatim}
\paragraph{What the ESS comparison does and does not show.}
The ESS experiments quantify autocorrelation costs for four
classical baselines under this artifact protocol.  The
parallel-tempering result (TODO-E2-MEAN-RATIO mean ratio) is the
most conservative comparator and the most appropriate headline
for any claim about residual quantum benefit.  The amortized
wall-clock comparison against exact inverse-CDF sampling
(\S IV-A) further shows that any structural advantage of the
quantum sampler does not translate into raw throughput in the
regime where the full distribution can be precomputed.  Together
these results define the scope of what the present submission can
and cannot demonstrate; they do not rule out a different conclusion
in regimes that we do not study, such as larger $n$ where
$O(2^n)$ preprocessing is infeasible, or noisy-hardware deployments.
\end{verbatim}
\end{itemize}

% -----------------------------------------------------------------------------
\subsection*{Patch 17 --- \S V-E (Method Selection Guidelines), revise}
% -----------------------------------------------------------------------------

\textbf{Replace the bulleted guidelines with:}

\begin{verbatim}
\subsection{Method Selection Guidelines}

\begin{itemize}
\item \textbf{Use exact inverse-CDF sampling} when $n$ is small
enough that $O(2^n)$ preprocessing is feasible.  This is the
strongest classical baseline in this regime and is the dominant
choice on raw wall-clock.

\item \textbf{Use amplitude-encoded quantum sampling} in the same
$n$-regime when (a) downstream workflows specifically require
i.i.d.\ samples drawn through a quantum measurement (e.g., for
quantum-statistical testing), (b) the preprocessing cost can be
amortized across many sampling tasks, or (c) integration with a
larger quantum pipeline is desired.  Do not select it on
wall-clock grounds against inverse-CDF.

\item \textbf{Use parallel tempering or HMC-with-discrete-relaxation}
when $n$ is too large for full enumeration but the distribution is
peaked enough that single-site Gibbs mixes slowly.

\item \textbf{Use junction-tree or message-passing methods} when the
graph is tree-structured or has small treewidth.

\item \textbf{Variational quantum compression (depth-3 hardware-efficient
ansatz)} is not a recommended sampling method based on our
results: at every $n \le 12$ we tested, mean-field or MPS produces
better full-distributional fidelity at matched or smaller
parameter budgets.
\end{itemize}
\end{verbatim}

\noindent\emph{Rationale:} Honest about where the methods studied here win
and lose.  Removes the ``Quantum provides measurable advantage when many
independent samples are required'' bullet from the original, since the new
inverse-CDF baseline shows this advantage does not survive when
preprocessing is amortized.

% -----------------------------------------------------------------------------
\subsection*{Patch 18 --- \S VI (Conclusion), full rewrite}
% -----------------------------------------------------------------------------

\begin{verbatim}
\section{Conclusion}

We presented a characterization of autocorrelation costs for four
classical sampling baselines on discrete MRFs, against an
amplitude-encoded quantum sampler that produces structurally i.i.d.\
output.  Across 60 instances and five graph families, the
Quantum/baseline ESS ratios decrease monotonically from 16.35
(single-site Gibbs) to 7.29 (block Gibbs) to 1.82 (tuned-block
Gibbs) to TODO-E2-MEAN (parallel tempering), making clear how much
of the apparent quantum sampling advantage in the literature is
attributable to baseline choice.  Once the $O(2^n)$ classical
preprocessing required by the amplitude encoder is amortized into
the wall-clock comparison, exact inverse-CDF sampling
TODO-E1-DOMINANCE-PHRASE the quantum sampler.  Together these
results define a precise scope within which the amplitude encoder
is a defensible sampling tool and outside of which it is not.

Secondary results: the amplitude-encoding pipeline reaches
$F \approx 1.000$ at $n = 8, 10, 12$ on simulator backends (a
sanity check on a Qiskit primitive); a multi-trial MPS scaling
study through $n = 40$ establishes a classical ceiling of
$F = 0.721 \pm 0.059$ at $\chi = 32$; and a shallow
hardware-efficient VQC underperforms matched-budget MPS at every
$n \le 12$ we tested.  Several diagnostic null results
(clique-aware entanglement, mean-field beating VQC on global fit)
are reported in the body and we believe are useful for the
community.

We do not claim quantum advantage in this work.  The most
plausible paths to one in this problem class --- full QCGM
Hamiltonian simulation, deeper or problem-aware variational
ansatze, or noisy-hardware benchmarks at scale --- are clearly
identified as future work.

All code, data, and artifacts are publicly available at
https://github.com/arulrhikm/QuantumDGM under an MIT license.
\end{verbatim}

\noindent\emph{Rationale:} Drops every overclaim in the original conclusion
and centers the contribution where the data actually supports it.

% -----------------------------------------------------------------------------
\subsection*{Patch 19 --- Bibliography fixes}
% -----------------------------------------------------------------------------

\begin{enumerate}
\item \textbf{Reference [7]:} Replace
\verb|art. 57, 2024| with \verb|art. 37, 2024|.  The correct full citation:

\begin{verbatim}
N. Piatkowski and C. Zoufal, ``Quantum circuits for discrete graphical
models,'' Quantum Machine Intelligence, vol. 6, art. 37, 2024.
DOI: 10.1007/s42484-024-00175-y.
\end{verbatim}

\item \textbf{Delete reference [21]} (Robert and Casella, \emph{Monte Carlo
Statistical Methods}) --- not cited anywhere in the body.

\item \textbf{Delete reference [22]} (Cowles and Carlin) --- not cited
anywhere in the body.

\item Renumber subsequent references accordingly.  In particular,
[23] (Brooks et al., MCMC handbook) becomes [21], and [24] (Lauritzen and
Spiegelhalter) becomes [22].  Update the in-text citations in \S V-B
(``Classical MCMC [23]'') and \S V-E (``junction-tree algorithm is
polynomial [24]'') accordingly.
\end{enumerate}

% -----------------------------------------------------------------------------
\subsection*{Patch 20 --- Table XII (status table), revise}
% -----------------------------------------------------------------------------

\textbf{Replace the existing Table XII with:}

\begin{verbatim}
\begin{table}[t]
\caption{Status of main experimental contributions in this revision.
``Complete'' = full multi-seed results in
\texttt{paper\_table\_metrics.json}.}
\label{tab:status}
\centering
\begin{tabular}{lll}
\toprule
Contribution & Status & Key metric / outcome \\
\midrule
ESS characterization, 4 baselines, 60 instances
& Complete (after E1/E2)
& 16.35 / 7.29 / 1.82 / TODO-PT \\
Amortized wall-clock vs.\ inverse-CDF
& Complete (after E1)
& TODO-E1-RATIO\,$\times$ \\
Amplitude-encoding pipeline verification ($n=8,10,12$)
& Complete
& $F \approx 1.000$ on BQ SV \\
MPS scaling ($n \le 40$, 3 seeds)
& Complete
& $\chi = 32$ gives $F = 0.721$ at $n = 40$ \\
VQC vs.\ MPS matched budget
& Complete (negative result)
& VQC below MPS at every $n \le 12$ \\
Clique-vs-linear entanglement
& Complete (null result)
& $\Delta F = -2.3 \times 10^{-4}$ \\
Mean-field vs.\ VQC global fidelity
& Complete (negative result for VQC)
& MF beats VQC at $n = 8, 10$ \\
Hardware deployment
& Out of scope
& No real-device results in this submission \\
Full QCGM (exponential-speedup path)
& Out of scope
& Identified as future work \\
\bottomrule
\end{tabular}
\end{table}
\end{verbatim}

\noindent\emph{Rationale:} Removes the ``Open / Partial'' rows that the
data no longer supports advertising as work-in-progress findings, and
explicitly marks the things that are simply not in this paper as
``Out of scope'' rather than ``Deferred.''

% =============================================================================
\section{Pre-Resubmission Checklist}
% =============================================================================

Before resubmitting, ensure the following:

\begin{enumerate}
\item Run E1 (inverse-CDF baseline) and write
      \verb|inverse_cdf_baseline.json|.
\item Run E2 (parallel tempering baseline) and write
      \verb|parallel_tempering_baseline.json|.
\item Update \verb|paper_table_metrics.json| to include the four-comparator
      summary and the amortized wall-clock numbers.
\item Apply Patches 1--20 above (text edits, no new experiments).
\item Replace every \verb|TODO-*| placeholder in the patched text with the
      corresponding number from the regenerated artifacts.
\item Drop the fidelity panel from Fig.~1 (keep runtime panel only).
\item Verify that no in-text citation refers to deleted references [21] or
      [22].
\item Re-run \verb|experiments/verify_paper_artifacts.py| to confirm every
      table number matches the JSON.
\item Final pass: search the paper for the words ``advantage,'' ``speedup,''
      and ``crossover.''  Each remaining occurrence must be honestly
      qualified by the relevant scope statement.
\end{enumerate}

\end{document}